\section{Cơ sở lý thuyết}

\subsection{Kiến trúc tổng quan (High-level design)}

Hệ thống được thiết kế theo mô hình pipeline streaming một chiều, trong đó dữ liệu chảy từ nguồn đến đích mà không cần lưu trữ trung gian lâu dài:

\begin{center}
\texttt{Producer → Kafka → Flink → Model → Kafka → Worker → PostgreSQL → BFF → Dashboard}
\end{center}

Các thành phần chính:

\begin{itemize}
    \item \textbf{Apache Kafka}: Message broker phân tán, đảm bảo thông lượng cao và độ bền dữ liệu. Hệ thống dùng 2 topic: \texttt{social\_media\_stream} (raw) và \texttt{final\_comments} (đã phân tích).
    \item \textbf{Apache Flink}: Stream processing engine xử lý dữ liệu với độ trễ mili-giây, thực hiện tiền xử lý văn bản.
    \item \textbf{RoBERTa Model}: Mô hình ngôn ngữ lớn thực hiện phân loại cảm xúc.
    \item \textbf{PostgreSQL}: Lưu trữ kết quả với schema quan hệ chuẩn hóa.
    \item \textbf{BFF + Dashboard}: Lớp API Express.js và giao diện React cung cấp khả năng trực quan hóa.
\end{itemize}

\subsection{Kiến trúc chi tiết (Low-level design)}

\subsubsection{Kafka Topics và Schema}

Hệ thống sử dụng 3 Kafka topic:

\begin{itemize}
    \item \texttt{social\_media\_stream}: Nhận raw review từ producer, message là JSON với các trường \texttt{text}, \texttt{timestamp}, \texttt{asin}, \texttt{star\_rating}.
    \item \texttt{processed\_comments}: Nhận output từ Flink sau khi làm sạch văn bản.
    \item \texttt{final\_comments}: Nhận output từ model sau khi phân tích cảm xúc, bổ sung \texttt{sentiment\_label} và \texttt{confidence\_score}.
\end{itemize}

\subsubsection{Apache Flink Job}

Flink job (\texttt{Main.java}) thực hiện:
\begin{enumerate}
    \item Đọc từ topic \texttt{social\_media\_stream} qua \texttt{KafkaSource}
    \item Deserialize JSON thành object \texttt{SocialMediaComment}
    \item Làm sạch văn bản: loại bỏ URL (\texttt{http://...}), @mention, \#hashtag bằng regex
    \item Serialize và ghi vào topic \texttt{processed\_comments} qua \texttt{KafkaSink}
\end{enumerate}

\subsubsection{Database Schema}

PostgreSQL lưu trữ theo 5 bảng:
\begin{itemize}
    \item \textbf{AmazonReview}: Thông tin gốc (review\_id, product\_asin, star\_rating, original\_text, event\_time)
    \item \textbf{SentimentResult}: Kết quả phân tích (review\_id FK, sentiment\_label, confidence\_score, model\_version\_id FK)
    \item \textbf{Product}: Danh sách sản phẩm theo ASIN
    \item \textbf{MLModel}: Thông tin phiên bản model
    \item \textbf{StreamTopic}: Metadata về Kafka topics
\end{itemize}

\subsubsection{BFF API Endpoints}

Express.js BFF (port 3001) cung cấp 4 endpoint:

\begin{table}[h]
\centering
\begin{tabular}{|l|l|l|}
\hline
\textbf{Endpoint} & \textbf{Params} & \textbf{Mô tả} \\
\hline
\texttt{GET /api/products} & --- & Danh sách ASIN \\
\hline
\texttt{GET /api/overview} & dateRange, productAsin, confidenceMin/Max & Thống kê tổng quan \\
\hline
\texttt{GET /api/trend} & dateRange, productAsin & Xu hướng theo ngày \\
\hline
\texttt{GET /api/reviews} & dateRange, productAsin, page, limit, confidenceMin/Max & Danh sách review \\
\hline
\end{tabular}
\end{table}

\subsection{Mô hình RoBERTa}

\subsubsection{Giới thiệu RoBERTa}

RoBERTa (Robustly Optimized BERT Pretraining Approach) là mô hình ngôn ngữ lớn được phát triển bởi Facebook AI, cải tiến từ BERT với các kỹ thuật:
\begin{itemize}
    \item Huấn luyện lâu hơn với batch size lớn hơn
    \item Loại bỏ next sentence prediction objective
    \item Huấn luyện trên tập dữ liệu lớn hơn (160GB text)
    \item Dynamic masking thay vì static masking
\end{itemize}

\subsubsection{Model được sử dụng}

Đồ án sử dụng checkpoint \texttt{cardiffnlp/twitter-roberta-base-sentiment} từ Hugging Face:
\begin{itemize}
    \item Fine-tuned trên 58 triệu tweet từ Twitter
    \item 3 nhãn đầu ra: \textbf{positive}, \textbf{neutral}, \textbf{negative}
    \item Input tối đa 512 token
    \item Tokenizer tương ứng: \texttt{AutoTokenizer} với \texttt{normalization=True}
\end{itemize}

\subsubsection{Quy trình inference}

\begin{enumerate}
    \item Đọc message từ Kafka topic \texttt{processed\_comments}
    \item Tokenize văn bản với \texttt{AutoTokenizer}
    \item Chạy forward pass qua mô hình, lấy logits
    \item Áp dụng softmax để ra xác suất 3 nhãn
    \item Chọn nhãn có xác suất cao nhất làm \texttt{sentiment\_label}
    \item Giá trị xác suất cao nhất là \texttt{confidence\_score}
    \item Publish kết quả vào topic \texttt{final\_comments}
\end{enumerate}

Để tối ưu hiệu suất, \texttt{model.py} xử lý theo batch nhỏ (default 8 messages/batch) và chạy inference trên CPU hoặc GPU tùy môi trường.
